% AGENTS: sketch of the Experimental setup section, drafted from the model code
% (src/bees/model.py), the experiment configs (configs/*.json), the experiment
% drivers (experiments/*.py), and the working report (report/report.md).
% Notation follows the Method section: B is the vertical-comb benefit,
% sigma_m the mutation rate, rho the sender--receiver coupling, D the maximum
% search distance, eta the worker-variation scale, kappa_max the maximum dance
% concentration. Numbers are the values actually used in the runs.

We study the model of \Cref{sec:method} in two stages that mirror the two
questions of \Cref{sec:intro}. The first stage keeps the comb horizontal and asks
under what ecological conditions the direct-pointing code is worth maintaining
(\Cref{sec:setup-horizontal}). The second stage releases the comb and the
transposition traits and asks under what conditions a population can migrate from
the direct to the gravity-referenced code as an extrinsic pressure tilts the comb
(\Cref{sec:setup-transition}). All experiments are run from configuration files
and fixed random seeds so that every reported run is reproducible.

\subsection{Fixed model parameters}
\label{sec:setup-fixed}
A core set of parameters is held fixed throughout, and only the ecological and
evolutionary parameters named below are varied. \Cref{tab:fixed-params} lists the
fixed values. The population contains $60$ colonies of $80$ workers each,
evolving for $120$ non-overlapping generations; each colony is evaluated over
$50$ independent foraging episodes of $12$ foraging attempts. Workers express
colony traits with variation scale $\eta = 0.08$ (\Cref{sec:traits}); dances are
emitted with maximum concentration $\kappa_{\max} = 14$ and production noise
$0.18$, and read with interpretation noise $0.12$ (\Cref{sec:signals}). All
lengths are in meters: the maximum search distance is $D = 6000$ m and food
patches lie at distances in $[750, d_{\max}]$ m. Food patches are physical disks
whose radius is drawn from a lognormal of median $150$ m and log-scale spread
$0.6$; capacity scales with patch area from a base of $6$ loads at the median
radius, and the number of patches per episode is Poisson (\Cref{sec:environment}).
Each foray draws its outbound distance from a Gamma of shape $2$ with the colony's
evolved mean. Foraging carries a per-distance travel cost, a dance cost
$c_{\mathrm{base}} + c_{\mathrm{cue}}\,b_i$ with $c_{\mathrm{base}} = 0$ and
$c_{\mathrm{cue}} = 0.02$, and an attention cost $c_{\mathrm{att}} = 0.01$
(\Cref{sec:foraging}). The sun is drawn from an arc of width $\pi$ centered on
$\pi/2$, and the comb orientation is treated as axially symmetric (period $\pi$).
Colonies start with little communicative ability: the initial directional bias is
drawn from $\mathcal{U}(0,0.15)$, receiver attention from $\mathcal{U}(0,0.25)$,
the foray distance from $\mathcal{U}(0.15D, 0.45D)$, the dance propensity from
$\mathcal{U}(0.8,1.0)$, and both transposition traits start at $0$, so early
foragers signal weakly, rarely attend to one another, and use only direct
pointing.

\begin{table}
  \centering
  \begin{tabular}{lll}
    \toprule
    Symbol / name & Value & Meaning \\
    \midrule
    colonies & $60$ & population size \\
    workers per colony & $80$ & ensemble size per colony \\
    generations & $120$ & evolutionary horizon (transition runs) \\
    episodes per colony & $50$ & foraging episodes averaged for fitness \\
    attempts per episode & $12$ & foraging attempts per episode \\
    $\eta$ & $0.08$ & worker-variation scale \\
    $\kappa_{\max}$ & $14$ & maximum dance concentration \\
    dance noise & $0.18$ & production noise (wrapped Gaussian) \\
    interpretation noise & $0.12$ & receiver decoding noise \\
    $D$ & $6000$ m & maximum search distance \\
    $d_{\min}$ & $750$ m & minimum food distance \\
    $\rho_r$ & $150$ m & median patch radius (lognormal) \\
    $\sigma_r$ & $0.6$ & patch-radius log-scale spread \\
    foray shape & $2$ & Gamma shape of per-foray distance \\
    $c_{\mathrm{base}},\,c_{\mathrm{cue}}$ & $0,\ 0.02$ & dance cost terms \\
    $c_{\mathrm{att}}$ & $0.01$ & attention cost \\
    $v$ & $1$ & food value \\
    $\mu_{\mathrm{sun}},\,\Delta_{\mathrm{sun}}$ & $\pi/2,\ \pi$ & sun arc center and width \\
    \bottomrule
  \end{tabular}
  \caption{Parameters held fixed across all experiments. Ecological parameters
    (food-site count, patch radius, capacity, maximum distance, travel cost) and
    evolutionary parameters ($B$, $\sigma_m$, $\rho$) are varied as described in
    the text.}
  \label{tab:fixed-params}
\end{table}

The parameters that are deliberately varied are summarized in
\Cref{tab:varied-params}: five ecological parameters (food-site count, patch
radius, capacity, maximum food distance, and travel cost) and three evolutionary
parameters ($B$, $\sigma_m$, $\rho$). The ranges given there are the union across
sub-experiments; the specific sweeps and grids used in each are described in the
sections below.

\begin{table}
  \centering
  \begin{tabular}{lll}
    \toprule
    Symbol / name & Range explored & Meaning \\
    \midrule
    \multicolumn{3}{l}{\emph{Ecological}} \\
    food-site count       & $1$--$8$          & Poisson mean number of patches per episode \\
    patch radius          & $60$--$360$ m     & median disk radius (lognormal) \\
    capacity              & $2$--$12$         & base loads a patch supplies at the median radius \\
    $d_{\max}$            & $3000$--$7500$ m  & maximum patch distance \\
    travel cost (per km)  & $0.010$--$0.060$ & cost charged per km of foray distance \\
    \midrule
    \multicolumn{3}{l}{\emph{Evolutionary}} \\
    $B$        & $0.10$--$0.60$ & vertical-comb benefit magnitude \\
    $\sigma_m$ & $0.04$--$0.14$ & mutation step standard deviation \\
    $\rho$     & $0.0$--$1.0$   & sender--receiver mutation correlation \\
    \bottomrule
  \end{tabular}
  \caption{Parameters varied across experiments. The \emph{Range explored} column
    gives the interval each parameter is sampled over in the Optuna search; the
    other sub-experiments use fixed grids that can extend beyond it. The ecological
    parameters and $B,\sigma_m,\rho$ are sampled continuously in the Optuna search
    over the intervals shown; the horizontal-comb experiment sweeps a grid of
    food-site count $\{1,2,3,4,6,8,12,16,24\}$ against median patch radius
    $\{15, 37.5, 75, 150, 300, 600\}$ m; and the evolutionary-interaction grid
    uses the discrete values $B \in \{0.10, 0.30, 0.45, 0.60\}$,
    $\sigma_m \in \{0.05, 0.07, 0.09, 0.11\}$, and
    $\rho \in \{0.0, 0.3, 0.6, 0.9\}$. Outside the search, capacity is $6$ at the
    median radius, $d_{\max} = 6000$ m, and the horizontal runs use
    $\sigma_m = 0.07$.}
  \label{tab:varied-params}
\end{table}

\subsection{Decoding the direct code: two variants}
\label{sec:setup-decode}
The direct code can be decoded in two ways on a tilted comb
(\Cref{sec:vertical}), and we run the whole transition pipeline under each. The
\emph{flatten} decode applies $M^\top$, introducing a tilt-dependent directional
bias; the \emph{unproject} decode applies $M^{-1}$, recovering the heading without
bias. The two pipelines are identical in every other respect, so comparing them
isolates the effect of the geometric decode assumption on whether and where the
transition occurs.

\subsection{Emergence of direct pointing (horizontal comb)}
\label{sec:setup-horizontal}
The first stage uses the horizontal-comb model of \Cref{sec:horizontal}: comb
tilt is fixed at $0$ and not allowed to evolve, so the gravity reference has zero
strength and the transposition traits are inert, leaving only the direct-pointing
code in play. Colonies evolve for $60$ generations under mutation rate
$\sigma_m = 0.07$, and we vary only the food distribution across $50$ replicate
seeds ($400$--$449$). The conditions form a full grid crossing the Poisson mean
patch count $\{1,2,3,4,6,8,12,16,24\}$ against the median patch radius
$\{15, 37.5, 75, 150, 300, 600\}$ m, so that the ecology ranges from single small
patches (essentially undiscoverable) to many large ones (easily found by random
search). All other ecological parameters are held at their fixed values
(\Cref{tab:fixed-params}).

The primary outcome is the \emph{recruitment advantage} (\Cref{sec:foraging}).
Because the follow decision is a randomized coin flip on receiver attention within
the same episodes, it is a near-randomized, contemporaneous estimate of what the
dance actually buys, and it distinguishes functioning communication from a
directional-bias trait that has merely drifted upward. We report it as a mean over
the final generations, alongside final-generation mean directional bias, foraging
success, and payoff.

\subsection{Evolution of the gravity code (vertical transition)}
\label{sec:setup-transition}
The second stage uses the full model of \Cref{sec:vertical}: comb tilt $\gamma$
and orientation $\phi$ are free to evolve, the transposition traits $t_s, t_r$ are
active, and the raw payoff is scaled by the vertical-comb benefit
$1 + B\gamma$. Populations always start horizontal, so a successful run must
evolve both a vertical comb \emph{and} a matched sender--receiver gravity code.
The transition is governed by the interaction of the ecology with three
evolutionary parameters: the benefit magnitude $B$, the mutation rate
$\sigma_m$, and the sender--receiver coupling $\rho$. Because the joint search
space is large, we locate viable regions by optimization and then probe their
robustness, in the sequence of sub-experiments below. Tuning and evaluation draw
from disjoint seed blocks: the search uses seeds $100$--$109$ and confirmation
$110$--$149$, while all held-out evaluation uses seeds $200$--$299$.

\paragraph{Parameter search.}
We search eight parameters with the Tree-structured Parzen Estimator sampler of
Optuna: food-site count, patch radius, capacity, maximum distance, and travel
cost, together with $B$, $\sigma_m$, and $\rho$ (food value fixed at $v = 1$). We
run $1024$ trials in total across sixteen parallel workers evaluated against a
single shared study, each worker seeded independently and drawing $64$ startup
random trials. Each trial evaluates a panel of $10$ seeds ($100$--$109$) and is
scored by
\begin{equation}
  \text{stable seed count} + \overline{\text{progress}} - \text{collapse count},
\end{equation}
where a seed is \emph{stable} if it ends with mean comb tilt $\ge 0.80$ and both
mean transposition traits $\ge 0.50$, \emph{collapsed} if its minimum
success rate over the run falls to $\le 0.02$, and \emph{progress} is a bounded
$[0,1]$ score measuring how close a non-stable seed came to the stable thresholds.
The stable-count term dominates while the progress term ranks near misses.

\paragraph{Confirmation and held-out validation.}
The strongest search trials are re-run on a disjoint confirmation panel of $40$
seeds ($110$--$149$), and the best candidates are then validated on $100$
held-out seeds ($200$--$299$) that were used in neither search nor confirmation.
Validation reports, per candidate, the fraction of stable seeds, the collapse
count, and final-generation means of foraging success, comb tilt, and the lower of
the two transposition traits.

\paragraph{One-parameter sensitivity.}
Around the strongest validated candidate we perturb each parameter in turn,
holding the others at their validated values, over the same $100$ held-out seeds
($200$--$299$), to identify which parameters the transition is most sensitive to.

\paragraph{Evolutionary-parameter interaction.}
To ask whether mutation parameters can compensate for a weaker architectural
benefit, we hold the validated ecology fixed and run a full-factorial grid over
the three evolutionary parameters, sweeping the benefit across the full search
range while re-centring the mutation parameters on the disk stable region,
$B \in \{0.10, 0.30, 0.45, 0.60\}$, $\sigma_m \in \{0.05, 0.07, 0.09, 0.11\}$, and
$\rho \in \{0.0, 0.3, 0.6, 0.9\}$, evaluating each of the $64$ cells over the same
$100$ held-out seeds ($200$--$299$).

\paragraph{Outcome measures.}
Across the transition experiments we summarize runs by the same thresholds used in
the search objective: the \emph{stable} fraction (vertical comb with a coordinated
gravity code), the \emph{gravity-reached} fraction (both transposition traits
crossing $0.50$ at any generation, whether or not verticality is retained), and
the \emph{collapse} fraction (foraging success falling to $\le 0.02$), alongside
final-generation trait and success means. Seed-bootstrap intervals over the
per-seed outcomes quantify sampling uncertainty in the reported fractions.
